\subsection{Invariant Distances and Adaptive Neighborhoods}
\begin{method} \label{meth:sl-invariant-and-tangent-distance} \hlt{Invariant Metrics and Tangent Distance}\\
Let $x,x'\in\R^p$, let $\mathcal{A}$ be a transformation family, and write $T_\alpha x$ for the transformed input.\\
For local transformations, let $V_x\in\R^{p\times r_x}$ span the tangent directions at $x$.
\begin{enumerate}[label=(\roman*)]
\setlength{\itemsep}{0pt}
\item \hlt{Orbit metric}: $d_{\mathrm{inv}}(x,x')=\inf\limits_{\alpha,\beta\in\mathcal{A}}\norm{T_\alpha x-T_\beta x'}_2$ compares the transformation orbits of $x$ and $x'$.
\item \hlt{Tangent metric}: $d_{\mathrm{tan}}(x,x')=\inf\limits_{a,b}\norm{x+V_xa-x'-V_{x'}b}_2$ compares tangent affine spaces.
\end{enumerate}
The tangent metric discounts small local transformations without making large transformations cost-free.\\
Data augmentation encodes similar invariance by storing transformed copies, which becomes costly when many transformation directions are needed.
\end{method}

\begin{figure}[H]
\centering
\begin{tikzpicture}[
  >=Stealth,
  every node/.style={font=\scriptsize},
  orbit/.style={green!65!black, line width=0.9pt},
  tangent/.style={red!75!black, line width=0.75pt},
  euclid/.style={blue!70!black, line width=0.7pt},
  point/.style={circle, fill=black, inner sep=1.45pt},
  guide/.style={gray!50, dashed, line width=0.45pt},
  labelbox/.style={fill=white, inner sep=1.2pt}
]
  \path[use as bounding box] (-0.25,0.35) rectangle (6.0,2.65);

  \coordinate (x) at (1.25,1.30);
  \coordinate (xp) at (4.75,1.30);

  \draw[orbit] (0.52,0.40)
    .. controls (0.72,0.78) and (0.96,1.02) .. (x)
    .. controls (1.52,1.56) and (1.60,2.20) .. (2.05,2.48);
  \draw[orbit] (3.95,2.48)
    .. controls (4.40,2.20) and (4.48,1.56) .. (xp)
    .. controls (5.04,1.02) and (5.28,0.78) .. (5.48,0.40);

  \node[point,label={[labelbox]below:$x_i$}] at (x) {};
  \node[point,label={[labelbox]below:$x_{i'}$}] at (xp) {};

  \draw[tangent] (0.70,0.78) -- (2.10,2.11);
  \draw[tangent] (3.90,2.11) -- (5.30,0.78);

  \draw[euclid] (x) -- (xp);
  \draw[tangent, <->] (1.85,1.87) -- (4.15,1.87);

  \draw[guide] (x) -- (1.85,1.87);
  \draw[guide] (xp) -- (4.15,1.87);

  \node[labelbox, text=green!55!black] at (1.35,2.45) {$\mathcal O_i$};
  \node[labelbox, text=green!55!black] at (4.65,2.45) {$\mathcal O_{i'}$};
  \node[labelbox, text=blue!70!black] at (3.00,1.12) {$d_E$};
  \node[labelbox, text=red!75!black] at (3.00,2.03) {$d_{\tan}$};
\end{tikzpicture}

\caption{Tangent distance uses \textcolor{red!75!black}{local tangents} instead of \textcolor{green!65!black}{orbits}; \textcolor{blue!70!black}{$d_E$} is ordinary distance.}
\label{fig:sl-tangent-distance-schematic}
\end{figure}

\begin{method} \label{meth:sl-discriminant-adaptive-nearest-neighbor} \hlt{Discriminant Adaptive Nearest Neighbors}\\
DANN replaces the fixed Euclidean metric by a query-dependent quadratic metric.\\
For a query $x_0$, choose a pilot index set $P(x_0)$, final neighbor count $h$, and $\epsilon>0$. From $\{(x_i,g_i):i\in P(x_0)\}$ estimate local class proportions $\pi_k$, class means $\bar{x}_k$, local mean $\bar{x}$, and pooled within-class covariance $\mathbf{W}\succ0$.
\begin{enumerate}[label=(\roman*)]
\setlength{\itemsep}{0pt}
\item Compute the local between-class covariance $\mathbf{B}=\sum\limits_{k=1}^K\pi_k(\bar{x}_k-\bar{x})(\bar{x}_k-\bar{x})^T$.
\item Form $\mathbf{B}^\star=\mathbf{W}^{-1/2}\mathbf{B}\mathbf{W}^{-1/2}$ and $\bm{\Sigma}_{x_0}=\mathbf{W}^{-1/2}(\mathbf{B}^\star+\epsilon I)\mathbf{W}^{-1/2}$.
\item With $D_{x_0}(x_i,x_0)=(x_i-x_0)^T\bm{\Sigma}_{x_0}(x_i-x_0)$, let $N_h(x_0)$ index $h$ smallest distances.\\
\item Set $\widehat{G}(x_0)\in\operatorname*{arg\,max}\limits_{g\in\{1,\ldots,K\}}\sum\limits_{i\in N_h(x_0)}\mathbf{1}_{\{g_i=g\}}$.
\end{enumerate}
Whitening by $\mathbf{W}$ removes local within-class scale.\\
$\mathbf{B}^\star+\epsilon I$ stretches neighborhoods along directions with weak local class separation, and $\epsilon$ prevents degeneracy.
\end{method}

\begin{figure}[H]
\centering
\begin{tikzpicture}[
  >=Stealth,
  every node/.style={font=\scriptsize},
  axis/.style={->, gray!70, line width=0.45pt},
  classa/.style={circle, draw=orange!85!black, fill=orange!12, inner sep=1.5pt},
  classb/.style={circle, draw=blue!75!black, fill=blue!10, inner sep=1.5pt},
  query/.style={circle, fill=black, inner sep=1.6pt},
  ordinary/.style={gray!65, dashed, line width=0.75pt},
  adapted/.style={green!55!black, line width=0.95pt},
  change/.style={red!75!black, line width=0.75pt, ->},
  labelbox/.style={fill=white, inner sep=1.1pt}
]
  \path[use as bounding box] (-0.45,-0.25) rectangle (6.25,3.65);

  \draw[axis] (0,0) -- (5.95,0) node[right] {$x_1$};
  \draw[axis] (0,0) -- (0,3.45) node[above] {$x_2$};
  \draw[gray!45, dashed, line width=0.55pt] (3.05,0.25) -- (3.05,3.15);

  \foreach \p in {(0.75,0.8),(1.15,1.55),(1.55,2.3),(2.15,0.95),(2.35,1.75),(2.55,2.75)}
    \node[classa] at \p {};
  \foreach \p in {(3.55,0.75),(3.75,1.55),(3.95,2.45),(4.55,1.0),(4.85,1.9),(5.25,2.8)}
    \node[classb] at \p {};

  \coordinate (q) at (3.05,1.85);
  \node[query,label={[labelbox]above right:$x_0$}] at (q) {};

  \draw[ordinary] (q) circle [radius=0.78];
  \draw[adapted] (q) ellipse [x radius=0.38, y radius=1.55];
  \draw[change] (2.42,1.85) -- (3.68,1.85);

  \node[labelbox, text=gray!65!black] at (2.25,2.75) {$B_E(x_0)$};
  \node[labelbox, text=green!45!black] at (3.72,3.10) {$B_D(x_0)$};
  \node[labelbox, text=red!75!black] at (3.05,1.48) {$\nabla\pi$};
\end{tikzpicture}

\caption{DANN replaces the \textcolor{gray!65!black}{Euclidean ball} by a \textcolor{green!55!black}{local ellipsoid}, narrow along the \textcolor{red!75!black}{class-change direction} and elongated where class probabilities vary slowly.}
\label{fig:sl-dann-adaptive-neighborhood}
\end{figure}

\begin{method} \label{meth:sl-global-dimension-reduction-nearest-neighbor} \hlt{Global Dimension Reduction for Nearest Neighbors}\\
Fix a target dimension $L\leq p$.  Let $\mathbf{B}_i$ be the local between-class matrix estimated at training point $x_i$.
\begin{enumerate}[label=(\roman*)]
\setlength{\itemsep}{0pt}
\item Compute the global class-separation matrix $\overline{\mathbf{B}}=N^{-1}\sum_{i=1}^N\mathbf{B}_i$.
\item Let $(\theta_\ell,e_\ell)_{\ell=1}^p$ be eigenpairs of $\overline{\mathbf{B}}$ with $\theta_1\geq\cdots\geq\theta_p$.
\item Project by $P_Lx=(e_1^Tx,\ldots,e_L^Tx)\in\R^L$ and run nearest-neighbor classification on $P_Lx_i$.
\end{enumerate}
The rank-$L$ summary $\overline{\mathbf{B}}_{[L]}=\sum_{\ell=1}^L\theta_\ell e_\ell e_\ell^T$ solves $\min_{\operatorname{rank}(M)=L}\sum_{i=1}^N\operatorname{trace}[(\mathbf{B}_i-M)^2]$.\\
The reduced space is selected from local class-separation information, not unsupervised variance alone.
\end{method}

\begin{remark} \label{rem:sl-nearest-neighbor-computation-and-storage} \hlt{Computation and Storage}\\
With $N$ stored observations in $\R^p$, direct nearest-neighbor search costs $O(Np)$ per query and stores $O(Np)$ coordinates.  Storage can be reduced by replacing the retained set $S\subset\{1,\ldots,N\}$.
\begin{enumerate}[label=(\roman*)]
\setlength{\itemsep}{0pt}
\item \hlt{Prototype compression}: store learned representatives rather than all training observations.
\item \hlt{Editing}: delete observations whose local neighborhoods have inconsistent labels.
\item \hlt{Condensing}: grow a small retained set by adding $i$ only if it misclassifies $(x_i,g_i)$.
\end{enumerate}
Fast search reduces query cost in favorable geometries; compression, editing, and condensing reduce storage by keeping representatives or boundary-relevant observations.
\end{remark}
